\subchapter{Bootloader - TF-A and U-Boot}{Objectives: Set up serial
  communication, compile and install the U-Boot bootloader, use basic
  U-Boot commands, set up TFTP communication with the development
  workstation.}

As the bootloader is the first piece of software executed by a
hardware platform, the installation procedure of the bootloader is
very specific to the hardware platform. There are usually two cases:

\begin{itemize}

\item The processor offers nothing to ease the installation of the
  bootloader, in which case the JTAG has to be used to initialize
  flash storage and write the bootloader code to flash. Detailed
  knowledge of the hardware is of course required to perform these
  operations.

\item The processor offers a monitor, implemented in ROM, and through
  which access to the memories is made easier.

\end{itemize}

The STM32MP2 SoC, falls into the second category. The monitor integrated in the
ROM reads the SD card to search for a valid bootloader
(the boot mode is actually configurable via a few input pins).
In case no bootloader is found, it will operate in a fallback mode, that will
allow to use an external tool to reflash some executable through USB.
Therefore, either by using an MMC/SD card or that fallback mode, we can start
up an STM32MP2-based board without having anything installed on it.
\section{Setup}

Go to the \code{$HOME/__SESSION_NAME__-labs/bootloader} directory.

\section{Setting up serial communication with the board}

Plug the USB-A or USB-C to USB-C cable in CN21 port also named USB PWR STLNK
of the Discovery board. It allows to power-on the board and also to expose
multiple debugging interfaces, including a serial interface.
When plugged in your computer, a serial port should appear, {\tt \hosttty}.

You can also see this device appear by looking at the output of
\code{sudo dmesg}.

To communicate with the board through the serial port, install a
serial communication program, such as \code{picocom}:

\bashcmd{$ sudo apt install picocom}

If you run {\tt ls -l \hosttty}, you can also see that only
\code{root} and users belonging to the \code{dialout} group have
read and write access to the serial console. Therefore, you need
to add your user to the \code{dialout} group:

\bashcmd{$ sudo adduser $USER dialout}

{\bf Important}: for the group change to be effective, you have to
reboot your computer (at least on Ubuntu 24.04) and log in again.
A workaround is to run \code{newgrp dialout}, but it is not global.
You have to run it in each terminal.

Run {\tt picocom -b 115200 \hosttty}, to start serial
communication on {\tt \hosttty}, with a baudrate of 115200.
If you wish to exit \code{picocom}, press \code{[Ctrl][a]} followed by
\code{[Ctrl][x]}.

Don't be surprised if you don't get anything on the serial console yet,
even if you reset the board. That's because the SoC has nothing to boot
on yet. We will prepare a micro SD card to boot on in the next paragraphs.

\section{TF-A and U-Boot relationship}

The boot process is done in two steps with the ROM monitor trying to
execute a first piece of software, called {\em fsbl}, from its
internal SRAM, that will initialize the DRAM, and a second program,
{\em ssbl} that will in turn load Linux and execute it.

In our case, {\em fsbl} is provided by TF-A BL2 and {\em
  ssbl} is provided by U-Boot.

TF-A BL2 is loading U-Boot from the Firmware Image Package ({\em
  FIP}), that will also contain the configuration for this second
part, a secure monitor (TF-A BL31) and a secure OS (OP-TEE BL32).
The {\em FIP} and the secure monitor are generated from TF-A sources,
so first we are going to build U-Boot and OP-TEE.

\section{U-Boot setup}

Download U-Boot:

\begin{bashinput}
$ git clone https://github.com/STMicroelectronics/u-boot.git
$ cd u-boot
$ git checkout v2023.10-stm32mp-r1
\end{bashinput}

Get an understanding of U-Boot's configuration and compilation steps
by reading the \code{README} file, and specifically the {\em Building
the Software} section.

Basically, you need to:

\begin{enumerate}

\item Specify the cross-compiler prefix
(the part before \code{gcc} in the cross-compiler executable name):
\bashcmd{$ export CROSS_COMPILE=aarch64-linux-}

\item Run \inlinebash{$ make stm32mp25_defconfig}

\item Now that you have a valid initial configuration, you can now
  run \inlinebash{$ make menuconfig} to further edit your bootloader features.

  \begin{itemize}
  \item In the \code{Device Drivers} $\rightarrow$ \code{Watchdog
      Timer Support} submenu, disable \code{IWDG watchdog
      driver for STM32 MP's family}, so that U-Boot doesn't start the
    watchdog.
  \end{itemize}

Install the following packages which should be needed to compile U-Boot for
your board:

\begin{bashinput}
$ sudo apt install libssl-dev device-tree-compiler swig \
       python3-dev python3-setuptools uuid-dev libgnutls28-dev
\end{bashinput}

\item Finally, run \bashcmd{make DEVICE_TREE=stm32mp257f-dk}
  which will build U-Boot
  \footnote{You can speed up the
  compiling by using the \code{-jX} option with \code{make}, where X
  is the number of parallel jobs used for compiling. Twice the
  number of CPU cores is a good value.}.

{\bf Note}: u-boot build may fail on your machine if you have a recent
version of python. Such issue is
\href{https://source.denx.de/u-boot/u-boot/-/commit/a63456b9191fae2fe49f4b121e025792022e3950}{already
fixed upstream}, but not in the version targeted for the training. To get the relevant fix,
you can cherry-pick the fix onto your local branch:

\code{git cherry-pick a63456b9191fae2fe49f4b121e025792022e3950}
\end{enumerate}

\section{OP-TEE setup}

Download OP-TEE:

\begin{bashinput}
$ git clone https://github.com/STMicroelectronics/optee_os.git
$ cd optee_os
$ git checkout 4.0.0-stm32mp-r1
\end{bashinput}

Several configuration parameters have to be passed to the Makefile:
\begin{itemize}
\item Specify the cross-compiler prefix (the part before gcc in the
      cross-compiler executable name), either using the environment
      variable:\inlinebash{$ export CROSS_COMPILE64=aarch64-linux-}, or just by
      adding it to the \code{make} commande line.
\item The architecture has to be selected: \code{ARCH=arm}
\item The STM32MP2 platform is selected too with \code{PLATFOM=stm32mp2}
\item Specify the device tree source file: \code{CFG_EMBED_DTB_SOURCE_FILE=stm32mp257f-dk.dts}
\item Specify log level to be printed with ERROR and INFO logs:
	\code{CFG_TEE_CORE_LOG_LEVEL=2}
\item Specify the location of all output files with \code{O=build}
\end{itemize}

We can now launch the build with a single command line:
\begin{bashinput}
$ make ARCH=arm CROSS_COMPILE64=aarch64-linux- PLATFORM=stm32mp2 \
  CFG_EMBED_DTB_SOURCE_FILE=stm32mp257f-dk.dts CFG_TEE_CORE_LOG_LEVEL=2 \
  O=build all
\end{bashinput}

At the end of the build, the important output files generated are
located in \code{build/core}. We will find there:

\begin{itemize}

\item \code{tee-header_v2.bin}, which is an header file used to load and
	identify OP-TEE.

\item \code{tee-pageable_v2.bin}, which is a binary that contains parts of the
  OP-TEE that are pageable, meaning they can be swapped in and out of memory
  as needed.

\item \code{tee-pager_v2.bin}, which is a binary used to manage pageable
  memory region for OP-TEE.

\end{itemize}

\section{TF-A setup}

Get TF-A sources from STMicroelectronics github:
\begin{bashinput}
$ cd ..
$ git clone https://github.com/STMicroelectronics/arm-trusted-firmware.git
$ cd trusted-firmware-a/
$ git checkout v2.10-stm32mp-r1
\end{bashinput}

Several configuration parameters have to be passed to the Makefile:
\begin{itemize}
\item Specify the cross-compiler prefix (the part before gcc in the
      cross-compiler executable name), either using the environment
      variable:\inlinebash{$ export CROSS_COMPILE=aarch64-linux-}, or just by
      adding it to the \code{make} commande line.
\item The architecture has to be selected: \code{ARCH=aarch64}, as
      well as the major version of Arm Architecture, here the Cortex A35 is
      an Armv8, so we need to use \code{ARM_ARCH_MAJOR=8}
\item The STM32MP2 platform is selected too with \code{PLAT=stm32mp2}
\item For this specific board, the device tree is generated and then
      needs to be specifed: \code{DTB_FILE_NAME=stm32mp257f-dk.dtb}
\item Specify the configuration of this firmware which is actually
      the Device Tree passed to U-Boot:
      \code{BL33_CFG=../u-boot/u-boot.dtb}
\item Specify that the fsbl will be located on the SD
      card with \code{STM32MP_SDMMC=1}.
\item Specify that the type of onboard RAM with \code{STM32MP_LPDDR4_TYPE=1}.
\item Specify the location of the BL33 (Boot loader stage 3-3):
      \code{BL33=../u-boot/u-boot-nodtb.bin}
\end{itemize}

We can now generate the \code{bl32}, \code{dtb}, and \code{fip} targets
with a single command line:
\begin{bashinput}
$ make ARM_ARCH_MAJOR=8 ARCH=aarch64 PLAT=stm32mp2 \
  DTB_FILE_NAME=stm32mp257f-dk.dtb \
  BL32=../optee_os/build/core/tee-header_v2.bin \
  BL32_EXTRA1=../optee_os/build/core/tee-pager_v2.bin \
  BL32_EXTRA2=../optee_os/build/core/tee-pageable_v2.bin \
  BL33=../u-boot/u-boot-nodtb.bin BL33_CFG=../u-boot/u-boot.dtb \
  STM32MP_SDMMC=1 STM32MP_LPDDR4_TYPE=1 fip all
\end{bashinput}

At the end of the build, the important output files generated are
located in \code{build/stm32mp2/release/}. We will find there:

\begin{itemize}

\item \code{tf-a-stm32mp257f-dk.stm32}, which is TF-A BL2, serving as
  our first stage bootloader

\item \code{fip.bin}, which is the FIP image, which itself includes
  U-Boot. This image will serve as the second stage bootloader.

\end{itemize}

\section{Flashing the bootloaders}

The ROM monitor will look for the first stage bootloader in a
partition named \code{fsbla1}. If it cannot find a valid bootloader in
this partition, it will then try to load it from a partition named
\code{fsla2}. This first stage bootloader (in our case the TF-A BL2)
will load the BL31 monitor (TF-A Secure monitor), the BL32 secure OS (OP-TEE)
and the second bootloader (U-Boot) from the Firmware Image
Package located in the partition named \code{fip}.
Then, the first stage bootloader jumps to secure monitor which then jumps to U-Boot.
Finally, U-Boot will store its environment in the default location.

So, as far as bootloaders are concerned, the SD card partitioning will
look like:

\begin{verbatim}
Number  Start   End      Size     File system  Name    Flags
 1      2048s   4095s    2048s                 fsbla1
 2      4096s   6143s    2048s                 fsbla2
 3      6144s   14335s   8191s                 fip
 4      14336s  131071s  135167s               bootfs
\end{verbatim}

On your workstation, plug in the SD card your instructor gave you. Type
the \code{sudo dmesg} command to see which device is used by your
workstation. In case the device is \code{/dev/mmcblk0}, you will see
something like

\begin{verbatim}
[46939.425299] mmc0: new high speed SDHC card at address 0007
[46939.427947] mmcblk0: mmc0:0007 SD16G 14.5 GiB
\end{verbatim}

The device file name may be different (such as \code{/dev/sdb}
if the card reader is connected to a USB bus (either internally
or using a USB card reader).

In the following instructions, we will assume that your SD card is
seen as \code{/dev/mmcblk0} by your PC workstation.

Type the \code{mount} command to check your currently mounted
partitions. If SD partitions are mounted, unmount them:

\bashcmd{$ sudo umount /dev/mmcblk0p*}

We will erase the existing partition table and partition contents
by simply zero-ing the first 128 MiB of the SD card:

\bashcmd{$ sudo dd if=/dev/zero of=/dev/mmcblk0 bs=1M count=128}

Now, let's use the \code{parted} command to create the partitions that
we are going to use:

\bashcmd{$ sudo parted /dev/mmcblk0}

The ROM monitor handles {\em GPT} partition tables, let's create one:

\begin{verbatim}
(parted) mklabel gpt
\end{verbatim}

Then, the 4 partitions are created with:
\begin{verbatim}
(parted) mkpart fsbla1 0% 4095s
(parted) mkpart fsbla2 4096s 6143s
(parted) mkpart fip 6144s 14335s
(parted) mkpart bootfs 14336s 135167s
\end{verbatim}

You can verify everything looks right with:

\begin{verbatim}
(parted) print
Model: SD SD16G (sd/mmc)
Disk /dev/mmcblk0: 15,5GB
Sector size (logical/physical): 512B/512B
Partition Table: gpt
Disk Flags: 

Number  Start   End     Size    File system  Name    Flags
 1      1049kB  2097kB  1049kB               fsbla1
 2      2097kB  3146kB  1049kB               fsbla2
 3      3146kB  7340kB  4194kB               fip
 4      7340kB  69,2MB  61,9MB               bootfs

(parted)
\end{verbatim}

Once done, quit:
\begin{verbatim}
(parted) quit
\end{verbatim}

{\em Note: \code{parted} is definitely not very user friendly compared
to other tools to manipulate partitions (such as \code{cfdisk}), but
that's the only tool which supports assigning names to GPT partitions.
In your projects, you could use \code{gparted}, which is a more
friendly graphical front-end on top of \code{parted}.}

Now, format the boot partition as an ext4 filesystem. This is where
U-Boot saves its environment:
\bashcmd{$ sudo mkfs.ext4 -L boot -O ^metadata_csum /dev/mmcblk0p4}

The \code{-O ^metadata_csum} option allows to create the filesystem
without enabling metadata checksums, which U-Boot doesn't seem to
support yet.

Now write the TF-A binary in both \code{fsbl} partitions:

\begin{bashinput}
$ sudo dd if=build/stm32mp2/release/tf-a-stm32mp257f-dk.stm32 of=/dev/mmcblk0p1 bs=1M conv=fdatasync
$ sudo dd if=build/stm32mp2/release/tf-a-stm32mp257f-dk.stm32 of=/dev/mmcblk0p2 bs=1M conv=fdatasync
\end{bashinput}

Then flash the {\em fip} partition with the Firmware Image Package
containing U-Boot, the BL32 monitor and their configuration (device tree):

\bashcmd{$ sudo dd if=build/stm32mp2/release/fip.bin of=/dev/mmcblk0p3 bs=1M conv=fdatasync}

\section{Testing the bootloaders}

Insert the SD card in the board slot. You can now power-up the board
by connecting the USB-C cable to the board, in CN6, \code{PWR_IN} and
to your PC at the other end. Check that it boots your new bootloaders.
You can verify this by checking the build dates:

\begin{verbatim}
NOTICE:  CPU: STM32MP257FAK Rev.Y
NOTICE:  Model: STMicroelectronics STM32MP257F-DK Discovery Board
NOTICE:  Board: MB1605 Var1.0 Rev.C-01
NOTICE:  BL2: v2.10-stm32mp2-r1.0(release):v2.10-stm32mp-r1(75f8c318)
NOTICE:  BL2: Built : 16:19:16, Mar 29 2025
NOTICE:  BL2: Booting BL31
NOTICE:  BL31: v2.10-stm32mp2-r1.0(release):v2.10-stm32mp-r1(75f8c318)
NOTICE:  BL31: Built : 16:19:08, Mar 29 2025
I/TC: Early console on UART#2
I/TC: 
I/TC: Embedded DTB found
I/TC: OP-TEE version: 4.0.0-stm32mp-r1 (gcc version 13.2.0 (crosstool-NG 1.26.0)) #1 Wed Mar 26 18:21:51 UTC 2025 aarch64
I/TC: WARNING: This OP-TEE configuration might be insecure!
I/TC: WARNING: Please check https://optee.readthedocs.io/en/latest/architecture/porting_guidelines.html
I/TC: Primary CPU initializing
I/TC: WARNING: All debug access are allowed
I/TC: Override the OTP 124: 0 to 0x18db6
I/TC: WARNING: Embeds insecure stm32mp_provisioning driver
I/TC: UART console (non-secure)
I/TC: PMIC STPMIC REFID:2.@ V1.1
I/TC: Platform stm32mp2: flavor PLATFORM_FLAVOR - DT stm32mp257f-dk.dts
I/TC: OP-TEE ST profile: secure_and_system_services
[    0.000000] SCP-firmware 2.13.0-intree-optee-os-4.0.0-stm32mp-r1
[    0.000000] 
[    0.000000] [FWK] Module initialization complete!
I/TC: Primary CPU switching to normal world boot
I/TC: Reserved shared memory is disabled
I/TC: Dynamic shared memory is enabled
I/TC: Normal World virtualization support is disabled
I/TC: Asynchronous notifications are enabled


U-Boot 2023.10-stm32mp-r1 (Mar 30 2025 - 16:23:50 +0200)

CPU: STM32MP257FAK Rev.Y
Model: STMicroelectronics STM32MP257F-DK Discovery Board
Board: stm32mp2 (st,stm32mp257f-dk)
Board: MB1605 Var1.0 Rev.C-01
DRAM:  4 GiB
optee optee: OP-TEE: revision 4.0 (d9c4df97)
I/TC: Reserved shared memory is disabled
I/TC: Dynamic shared memory is enabled
I/TC: Normal World virtualization support is disabled
I/TC: Asynchronous notifications are enabled
Core:  432 devices, 41 uclasses, devicetree: board
WDT:   Started watchdog with servicing every 1000ms (32s timeout)
NAND:  0 MiB
MMC:   STM32 SD/MMC: 0, STM32 SD/MMC: 1
Loading Environment from MMC... *** Warning - bad CRC, using default environment

In:    serial
Out:   serial
Err:   serial
Net:   eth0: eth1@482c0000
FWU metadata read failed
No EFI system partition
No EFI system partition
Failed to persist EFI variables
Hit any key to stop autoboot:  0 
Boot over mmc0!
Saving Environment to MMC... Writing to MMC(0)... OK
switch to partitions #0, OK
mmc0 is current device
STM32MP>
\end{verbatim}

In U-Boot, type the \code{help} command, and explore the few commands
available.

\subsection{Adding a new command to the U-Boot shell}

Check whether the \code{config} command is available. This command
allows to dump the configuration settings U-Boot was compiled from.

If it's not, go back to U-Boot's configuration and enable it.

Re-run the build of U-Boot, and then re-run the build of TF-A so that
a new version of the \code{fip.bin} with the updated U-Boot is
generated.

Update the \code{fip} partition on the SD card with the new
\code{fip.bin} image and test that the command is now available and
works as expected.

\subsection{Playing with the U-Boot environment}

Display the U-Boot environment using \code{printenv}.

Set a new U-Boot variable \code{foo} to a value of your choice, using
\code{setenv}, and verify it has been set. Reset the board, and check
if \code{foo} is still defined: it should not.

Now repeat this process, but before resetting the board, use
\code{saveenv}. After the reset, check the \code{foo} variable is
still defined.

\section{Setting up networking}

The next step is to configure U-boot and your workstation to let your
board download files, such as the kernel image and Device Tree Binary
(DTB), using the TFTP protocol through a network connection.

With a network cable, connect the Ethernet port of
your board to the one of your computer. If your computer already has a
wired connection to the network, your instructor will provide you with
a USB Ethernet adapter. A new network interface should appear on your
Linux system.

\subsection{Network configuration on the target}

Let's configure networking in U-Boot:

\begin{itemize}
  \item \code{ipaddr}: IP address of the board
  \item \code{serverip}: IP address of the PC host
\end{itemize}

\begin{ubootinput}
=> setenv ipaddr 192.168.0.100
=> setenv serverip 192.168.0.1
\end{ubootinput}

Of course, make sure that this address belongs to a separate network
segment from the one of the main company network.

To make these settings permanent, save the environment:

\begin{ubootinput}
=> saveenv
\end{ubootinput}

\subsection{Network configuration on the PC host}

To configure your network interface on the workstation side, we need
to know the name of the network interface connected to your board.

Find the name of this interface by typing:
\bashcmd{=> ip a}

The network interface name is likely to be
\code{enxxx}\footnote{Following the {\em Predictable Network Interface
Names} convention:
\url{https://www.freedesktop.org/wiki/Software/systemd/PredictableNetworkInterfaceNames/}}.
If you have a pluggable Ethernet device, it's easy to identify as it's
the one that shows up after pluging in the device.

Then, instead of configuring the host IP address from NetworkManager's
graphical interface, let's do it through its command line interface,
which is so much easier to use:

\bashcmd{$ nmcli con add type ethernet ifname en... ip4 192.168.0.1/24}

\section{Setting up the TFTP server}

Let's install a TFTP server on your development workstation:

\begin{verbatim}
sudo apt install tftpd-hpa
\end{verbatim}

You can then test the TFTP connection. First, put a small text file in
the directory exported through TFTP on your development
workstation. Then, from U-Boot, do:

\begin{ubootinput}
=> tftp %\zimageboardaddr% textfile.txt
\end{ubootinput}

The \code{tftp} command should have downloaded the
\code{textfile.txt} file from your development workstation into
the board's memory at location {\tt \zimageboardaddr}\footnote{
This location is part of the board DRAM. If you want
to check where this value comes from, you can check the SoC
datasheet at
\url{https://www.st.com/resource/en/reference_manual/rm0457-stm32mp2325xx-advanced-armbased-3264bit-mpus-stmicroelectronics.pdf}.
It's a big document (more than 4,000 pages). In this document, look
for \code{Memory organization} and you will find the SoC memory map.
You will see that the address range for the memory controller
({\em DDRC})
starts at the address we are looking for.
You can also try with other values in the RAM address range.}.

You can verify that the download was successful by dumping the
contents of the memory:

\begin{ubootinput}
=> md %\zimageboardaddr%
\end{ubootinput}

We will see in the next labs how to use U-Boot to download, flash and
boot a kernel.

\section{Rescue binaries}

If you have trouble generating binaries that work properly, or later
make a mistake that causes you to lose your bootloader binaries, you
will find working versions under \code{data/} in the current lab
directory.
